\documentclass[12pt]{article}
\usepackage[margin=1in,a4paper]{geometry}
\usepackage{amsmath,amssymb,fullpage,xcolor,mhchem,mdframed}
\usepackage{hyperref}
\DeclareMathOperator\Tr{Tr}
\begin{document}
\title{Chapter 40: Spin Sums}
\date{\today}
\maketitle

\begin{mdframed}[backgroundcolor=blue!5]
\textbf{Chapter 40 bridges the gap between Feynman rules (Ch.39) and
observable cross sections.} Every experimental prediction requires
summing over final spins and averaging over initial spins. This chapter
shows how to do that efficiently using trace theorems.
\end{mdframed}

\section{Spin-Averaging Factors}\label{sec:averaging}
For an unpolarized initial state:
\begin{itemize}
\item Fermion (spin-$\frac12$): average over 2 spin states $\to \frac12$
\item Photon/gluon: average over 2 physical polarizations $\to \frac12$ (QED/QCD gauge)
\end{itemize}
For $e^+e^-\to\mu^+\mu^-$ at high energy ($s\gg m_\mu^2$):
\begin{equation}
\frac{d\sigma}{d\Omega}
= \frac{\alpha^2}{4s}\left[1+\cos^2\theta\right] .
\tag{40.$\star$}
\end{equation}

\section{Trace Method}\label{sec:trace_method}
Squared amplitudes with closed fermion loops become traces:
\begin{mdframed}[backgroundcolor=green!5]
\begin{equation}
\frac{1}{4}\sum_{\rm spins}|M|^2
= \frac{1}{4}\,\Tr\bigl[(\p1\!\!\!/+m)\gamma^\mu(\p2\!\!\!/+m)\gamma^\nu\bigr]\;
\Tr\bigl[(\p3\!\!\!/+m')\gamma_\mu(\p4\!\!\!/+m')\gamma_\nu\bigr] .
\tag{40.1}
\end{equation}
\end{mdframed}
Each closed fermion loop contributes one trace.

\section{Trace Theorems}\label{sec:trace_theorems}
Fundamental identities ($\eta=\mathrm{diag}(-1,+1,+1,+1)$):
\begin{align}
\Tr[\mathbf{1}] &= 4, &
\Tr[\gamma^\mu\gamma^\nu] &= 4\eta^{\mu\nu}, \\[2mm]
\Tr[\gamma^\mu\gamma^\nu\gamma^\rho\gamma^\sigma]
&= 4\bigl(\eta^{\mu\nu}\eta^{\rho\sigma}
          -\eta^{\mu\rho}\eta^{\nu\sigma}
          +\eta^{\mu\sigma}\eta^{\nu\rho}\bigr) .
\end{align}

Contracted trace identity (used in every QED cross-section):
\begin{equation}
\Tr[(\p1\!\!\!/+m)\gamma^\mu(\p2\!\!\!/+m)\gamma^\nu]
= 2\Bigl[p_1^\mu p_2^\nu + p_1^\nu p_2^\mu
         -\eta^{\mu\nu}(p_1\cdot p_2-m^2)\Bigr] .
\tag{40.2}
\end{equation}

\section{Helicity vs. Trace Method}\label{sec:vs_helicity}
\begin{center}
\begin{tabular}{|l|p{0.35\linewidth}|p{0.35\linewidth}|}
\hline
Feature & Trace Method & Helicity Method (Ch.42) \\
\hline
Output & Scalar products $p_i\cdot p_j$ & $\langle ij\rangle$, $[ij]$ \\
Best for & Loops, fermions & Gluon tree amplitudes \\
Result size & $O(n!)$ terms & Single expression for MHV \\
\hline
\end{tabular}
\end{center}

\section{Klein-Nishina (Compton Scattering)}\label{sec:compton}
\begin{mdframed}[backgroundcolor=red!5]
\begin{equation}
\frac{d\sigma}{d\Omega}
= \frac{\alpha^2}{2m_e^2}
\left(\frac{\omega'}{\omega}\right)^2
\left[\frac{\omega'}{\omega}+\frac{\omega}{\omega'}
-\sin^2\theta\cos^2\phi\right] .
\tag{40.$\star\star$}
\end{equation}
\end{mdframed}
In the Thomson limit $\omega\gg m_e$: $\sigma\to\sigma_T=\frac{8\pi r_e^2}{3}$.

\section{Significance for MHV Research}
Ch.40 is the last chapter before Ch.42 that uses only Lorentz-covariant
methods.  The trace theorems here are replaced by Schouten identities in the
helicity formalism: the same algebraic structure, but packaged in spinor
brackets that make helicity selection rules transparent.

\end{document}
